% Originally: content/week-08.tex, \subsection{Fubini's theorem} --- promoted to a section

\section{Fubini's theorem}
\label{sec:fubini}

\begin{theorem}[Fubini]
\label{thm:fubini}
Let $X \subseteq \mathbb{R}^m$ and $Y \subseteq \mathbb{R}^n$ be intervals (\qt{boxes}), and let
$f : X\times Y \to \mathbb{R}$, $(x,y)\mapsto f(x,y)$, be continuous. Then
\[\int_{X\times Y} f(x,y)\,dx\,dy
  = \int_X\left(\int_Y f(x,y)\,dy\right)dx
  = \int_Y\left(\int_X f(x,y)\,dx\right)dy,\]
where we look at the expression in brackets as a function, e.g.\
$x \mapsto \int_Y f(x,y)\,dy$.
\end{theorem}

% Source: Diego Torres Tejeda/Notes - 27.04 - Diego Torres Tejeda.pdf, p. 1
% Generator: Gemini 3.6 Flash (High)
\begin{proposition}[Cavalieri's principle]
\label{prop:cavalieri_principle}
Let $A, B \subseteq \mathbb{R}^{n+1}$ be Jordan measurable sets. For each height $t \in \mathbb{R}$, denote by $A_t := \{x \in \mathbb{R}^n \mid (x,t) \in A\}$ and $B_t := \{x \in \mathbb{R}^n \mid (x,t) \in B\}$ the $n$-dimensional cross-sections at $t$. If the cross-sectional volumes agree for all $t \in \mathbb{R}$, i.e.\
\[
\mu_n(A_t) = \mu_n(B_t) \quad \forall t \in \mathbb{R},
\]
then the total $(n+1)$-dimensional volumes are equal: $\mu_{n+1}(A) = \mu_{n+1}(B)$.
\end{proposition}

\begin{proof}
By Fubini's Theorem (\cref{thm:fubini}), integrating the indicator function over slices gives:
\[
\mu_{n+1}(A) = \int_{\mathbb{R}^{n+1}} \chi_A(x,t) \,dx\,dt = \int_{\mathbb{R}} \left( \int_{\mathbb{R}^n} \chi_{A_t}(x) \,dx \right) dt = \int_{\mathbb{R}} \mu_n(A_t) \,dt.
\]
Replacing $\mu_n(A_t)$ with $\mu_n(B_t)$ yields $\mu_{n+1}(B)$.
\end{proof}

\begin{theorem}[Volume of a cone over a measurable base]
\label{thm:cone_volume_formula}
Let $\Omega \subseteq \mathbb{R}^n$ be a bounded, Jordan measurable domain. The \newterm{cone over $\Omega$} \germanfor{Kegel \"uber $\Omega$} with vertex at $(0,1) \in \mathbb{R}^n \times \mathbb{R}$ and base $\Omega \times \{0\}$ is defined by
\[
C\Omega := \big\{(x, t) \in \mathbb{R}^n \times [0,1] \mid x \in (1-t)\Omega\big\}.
\]
Then $C\Omega$ is Jordan measurable in $\mathbb{R}^{n+1}$, and its volume is given by
\[
\mu_{n+1}(C\Omega) = \frac{\mu_n(\Omega)}{n+1}.
\]
\end{theorem}

\begin{proof}
For any $t \in [0,1]$, the cross-section of $C\Omega$ at height $t$ is $(C\Omega)_t = (1-t)\Omega \subseteq \mathbb{R}^n$. By the scaling property of $n$-dimensional Jordan measure, $\mu_n\big((1-t)\Omega\big) = (1-t)^n \mu_n(\Omega)$. Applying Cavalieri's Principle (\cref{prop:cavalieri_principle}):
\[
\mu_{n+1}(C\Omega) = \int_0^1 \mu_n\big((1-t)\Omega\big) \,dt = \mu_n(\Omega) \int_0^1 (1-t)^n \,dt = \mu_n(\Omega) \left[ -\frac{(1-t)^{n+1}}{n+1} \right]_0^1 = \frac{\mu_n(\Omega)}{n+1}.
\]
\end{proof}

\begin{center}
% Generator: Gemini 3.6 Flash (High) --- FIG-CONE: Cavalieri's principle for the volume of a cone over a domain
\begin{tikzpicture}[>=Latex, scale=1.0]

  % Axes
  \draw[->, thick] (-0.5,0) -- (3.2,0) node[right, font=\footnotesize] {$\mathbb{R}^n$};
  \draw[->, thick] (0,-0.3) -- (0,2.6) node[above, font=\footnotesize] {$t \in \mathbb{R}$};

  % Vertex at (0, 1.8)
  \fill[red!80!black] (0, 1.8) circle (1.8pt) node[left, font=\footnotesize] {$(0,1)$};

  % Base domain Omega at t=0
  \draw[very thick, orange!90!black] (0.8,0) -- (2.6,0);
  \fill[orange!90!black] (0.8,0) circle (1.5pt);
  \fill[orange!90!black] (2.6,0) circle (1.5pt);
  \node[orange!90!black, font=\footnotesize, below] at (1.7, -0.05) {$\Omega$};

  % Cone sides
  \draw[blue!75!black, thick] (0, 1.8) -- (0.8,0);
  \draw[blue!75!black, thick] (0, 1.8) -- (2.6,0);

  % Shaded region for C\Omega
  \fill[blue!10, opacity=0.7] (0, 1.8) -- (0.8,0) -- (2.6,0) -- cycle;
  \node[blue!80!black, font=\small] at (1.1, 0.7) {$C\Omega$};

  % Cross-section at height t
  \draw[dashed, gray!70!black] (-0.3, 1.0) -- (2.8, 1.0);
  \draw[very thick, green!50!black] (0.355, 1.0) -- (1.155, 1.0);
  \fill[green!50!black] (0.355, 1.0) circle (1.4pt);
  \fill[green!50!black] (1.155, 1.0) circle (1.4pt);
  \node[green!50!black, font=\scriptsize, above] at (0.75, 1.05) {$(1-t)\Omega$};

  \node[font=\scriptsize, left] at (0, 1.0) {$t$};
  \node[font=\scriptsize, left] at (0, 0.0) {$0$};

\end{tikzpicture}
\end{center}

% Source: Corsin Nick/Class Notes/Week 8.pdf, p. 8
\begin{example}[Fubini on the unit square]
\label{ex:fubini_unit_square}
\[
\begin{aligned}
\int_{[0,1]^2} x^{y+1}\,dx\,dy
  &= \int_0^1 dy \int_0^1 dx\ x^{y+1}
   = \int_0^1 dy\ \left[\frac{x^{y+2}}{y+2}\right]_0^1 \\
  &= \int_0^1 dy\ \frac{1}{y+2}
   = \log|y+2|\,\Big|_0^1
   = \log\!\left(\tfrac32\right).
\end{aligned}
\]
\end{example}

\begin{ainote}
Corsin writes the first line as $\int_0^1 dx \int_0^1 dy$, but the antiderivative taken on the
next line is in $x$, and it produces $x^{y+2}/(y+2)$, evaluated from $0$ to $1$. The
differentials are therefore swapped relative to the computation actually performed. The order is corrected
above and the result is unaffected.

The example is well chosen: integrating $x^{y+1}$ in $y$ first would mean integrating an
exponential in the exponent, which is far worse. Fubini's real use is that you may pick the
easier order, and here only one of the two is comfortable.
\end{ainote}

% Supplement: Sascha Brack/Class Notes/Week_08_Notes_Friday.pdf, pp. 4--5
\begin{example}[Changing the order of integration]
\label{ex:change_order_integration}
Consider the iterated integral
\[\int_0^2 \int_{y^2}^{2\sqrt{2y}} f(x,y) \,dx\,dy.\]
Let us rewrite it with the integration order reversed (first $y$, then $x$).

The region is defined by $0 \leq y \leq 2$ and $y^2 \leq x \leq 2\sqrt{2y}$.
Notice that for $y \in [0,2]$, the lower bound is $x \geq 0$ and the upper bound is $x \leq 2\sqrt{4} = 4$, so $x$ ranges from $0$ to $4$.
We now solve the inequalities for $y$:
$x \leq 2\sqrt{2y} \implies x^2 \leq 8y \implies y \geq \frac{1}{8}x^2$.
$x \geq y^2 \implies y \leq \sqrt{x}$ (since $x,y \geq 0$).
Thus, the new bounds are $0 \leq x \leq 4$ and $\frac{1}{8}x^2 \leq y \leq \sqrt{x}$. The reversed integral is:
\[\int_0^4 \int_{\frac{1}{8}x^2}^{\sqrt{x}} f(x,y) \,dy\,dx.\]
\exinfo{The same region, with the same two bounds, is Problem 14 of the Analysis II examination of
August 2025 (Prof.\ Laura Kobel-Keller), a box problem where only the reversed integral was
graded.}
\end{example}

% Generator: Claude Opus 5 (High) --- FIG-20-01, the two slicings of one region.
% Requested in feedback/brainstorming_improvements.md, which names swapping bounds over a
% non-rectangular domain as a major pain point and asks for exactly this picture. Drawn against
% the region of ex:change_order_integration rather than the brainstorm's triangle, because a
% triangle makes the two orders look symmetric and this region does not: the curved boundaries
% are what force the algebra in the example above.
%
% Arithmetic check. The two boundary curves are the SAME two curves, read two ways:
%   x = y^2        <=>  y = sqrt(x)      (upper boundary)
%   x = 2 sqrt(2y) <=>  y = x^2/8        (lower boundary)
% They meet where sqrt(x) = x^2/8, i.e. x^{3/2} = 8, x = 4, y = 2 -- and at the origin.
% Both panels must therefore show the SAME shaded region; only the slice differs.
%   LEFT (dx dy, outer y): horizontal slice at y = 1.2, running x = y^2 = 1.44
%       to x = 2 sqrt(2.4) = 3.0984.
%   RIGHT (dy dx, outer x): vertical slice at x = 2.2, running y = 2.2^2/8 = 0.605
%       to y = sqrt(2.2) = 1.4832.
% Both slice endpoints lie on the boundary curves by construction, which is the property the
% figure asserts and the one worth re-checking if the coordinates are ever touched.
\begin{center}
\begin{tikzpicture}[x=1.05cm, y=1.05cm, >=Latex,
    bdry/.style={blue!70!black, very thick},
    slice/.style={BrickRed, very thick, <->}]

  % ---------------- left panel: horizontal slices, the order as given ----------------
  \begin{scope}
    \fill[blue!8]
      plot[domain=0:4, samples=60] ({\x}, {sqrt(\x)}) --
      plot[domain=4:0, samples=60] ({\x}, {(\x)*(\x)/8}) -- cycle;

    \draw[->, gray!70] (-0.25,0) -- (4.9,0) node[right, font=\footnotesize, text=black] {$x$};
    \draw[->, gray!70] (0,-0.25) -- (0,2.7) node[above, font=\footnotesize, text=black] {$y$};

    \draw[bdry] plot[domain=0:4, samples=60] ({\x}, {sqrt(\x)});
    \draw[bdry] plot[domain=0:4, samples=60] ({\x}, {(\x)*(\x)/8});
    % Placed at x = 3.4 on the upper curve (y = sqrt(3.4) = 1.844) rather than near x = 1: the
    % dashed guide at y = 1.2 runs across the left half of this panel and the label collided
    % with it there.
    \node[above left, font=\scriptsize, blue!70!black, inner sep=1pt] at (3.4,1.844) {$x=y^2$};
    \node[below right, font=\scriptsize, blue!70!black, inner sep=1pt] at (3.05,1.16)
      {$x=2\sqrt{2y}$};

    \draw[slice] (1.44,1.2) -- (3.0984,1.2);
    \draw[gray!60, dashed] (0,1.2) -- (1.44,1.2);
    \node[left, font=\scriptsize] at (-0.1,1.2) {$y$};

    \fill (4,2) circle (1.4pt);
    \node[right, font=\scriptsize, inner sep=2pt] at (4.05,2.1) {$(4,2)$};

    \node[font=\scriptsize, BrickRed] (Lslab) at (1.05,2.5) {$y^2 \le x \le 2\sqrt{2y}$};
    \draw[BrickRed, thin] (Lslab) -- (1.9,1.28);
    \node[font=\footnotesize, align=center] at (2.0,-0.95)
      {$\displaystyle\int_0^2\!\!\int_{y^2}^{2\sqrt{2y}} f \,dx\,dy$};
    \node[font=\scriptsize, gray!50!black] at (2.0,-1.75) {fix $y$, sweep $x$};
  \end{scope}

  % ---------------- right panel: vertical slices, the order after swapping ------------
  \begin{scope}[shift={(6.6,0)}]
    \fill[blue!8]
      plot[domain=0:4, samples=60] ({\x}, {sqrt(\x)}) --
      plot[domain=4:0, samples=60] ({\x}, {(\x)*(\x)/8}) -- cycle;

    \draw[->, gray!70] (-0.25,0) -- (4.9,0) node[right, font=\footnotesize, text=black] {$x$};
    \draw[->, gray!70] (0,-0.25) -- (0,2.7) node[above, font=\footnotesize, text=black] {$y$};

    \draw[bdry] plot[domain=0:4, samples=60] ({\x}, {sqrt(\x)});
    \draw[bdry] plot[domain=0:4, samples=60] ({\x}, {(\x)*(\x)/8});
    \node[above left, font=\scriptsize, blue!70!black, inner sep=1pt] at (1.15,1.07)
      {$y=\sqrt{x}$};
    \node[below right, font=\scriptsize, blue!70!black, inner sep=1pt] at (3.05,1.16)
      {$y=x^2/8$};

    \draw[slice] (2.2,0.605) -- (2.2,1.4832);
    \draw[gray!60, dashed] (2.2,0) -- (2.2,0.605);
    \node[below, font=\scriptsize] at (2.2,-0.08) {$x$};

    \fill (4,2) circle (1.4pt);
    \node[right, font=\scriptsize, inner sep=2pt] at (4.05,2.1) {$(4,2)$};

    \node[font=\scriptsize, BrickRed] (Rslab) at (1.15,2.5) {$\tfrac{x^2}{8} \le y \le \sqrt{x}$};
    \draw[BrickRed, thin] (Rslab) -- (2.2,1.55);
    \node[font=\footnotesize, align=center] at (2.0,-0.95)
      {$\displaystyle\int_0^4\!\!\int_{x^2/8}^{\sqrt{x}} f \,dy\,dx$};
    \node[font=\scriptsize, gray!50!black] at (2.0,-1.75) {fix $x$, sweep $y$};
  \end{scope}
\end{tikzpicture}
\end{center}

\begin{remark}[Reading the swap off the picture]
\label{rem:fubini_swap_picture}
The region is the same in both panels; only the direction of the slice changes, and with it which
variable is the outer one. That is the whole of Fubini's theorem for a non-rectangular domain, and
the two mechanical rules follow from looking:

\begin{itemize}
  \item \textbf{The outer limits are constants, and they are the shadow of the region on the outer
    axis.} On the left the outer variable is $y$ and the region's shadow on the $y$-axis is
    $[0,2]$; on the right the outer variable is $x$ and its shadow is $[0,4]$. The two are
    different numbers, which is the most common slip.
  \item \textbf{The inner limits are functions of the outer variable, read off where the slice
    enters and leaves the region.} Entering and leaving happens on the \emph{same two curves} in
    both panels --- there is only one pair of boundary curves here. What changes is which of
    $x = y^2$ and $y = \sqrt{x}$ is the useful way to write the first of them.
\end{itemize}

The reason the algebra in the example above is not merely a relabelling is the second bullet: each
boundary has to be re-solved for the new inner variable. A region whose boundary cannot be solved
for one of the variables cannot be swapped at all without first splitting it, and that is the
situation \cref{ex:change_of_variables_domain} runs into.
\end{remark}

\begin{example}[Surface parametrization and Fubini]
\label{ex:surface_parametrization_integral}
Let $f(u,v) := (u\cos v, u\sin v, u^2)\transp \in \mathbb{R}^3$ for $(u,v) \in \mathbb{R}^2$. Write the third coordinate $z(u,v)$ as a function of the first two, $\tilde{z}(x,y)$, and compute $\int_0^y \int_0^x \tilde{z}(x',y') \,dx'\,dy'$.

We have $x(u,v) = u\cos v$ and $y(u,v) = u\sin v$. Then $x^2 + y^2 = u^2\cos^2 v + u^2\sin^2 v = u^2 = z(u,v)$.
So, we define $\tilde{z}(x,y) := x^2 + y^2$.
The integral is:
\begin{align}
\int_0^y \int_0^x (x'^2 + y'^2) \,dx'\,dy'
&= \int_0^y \left[ \frac{x'^3}{3} + x'y'^2 \right]_{x'=0}^{x'=x} \,dy' \nonumber \\
&= \int_0^y \left( \frac{x^3}{3} + x y'^2 \right) \,dy' \nonumber \\
&= \left[ \frac{x^3 y'}{3} + x \frac{y'^3}{3} \right]_{y'=0}^{y'=y} \nonumber \\
&= \frac{xy}{3}(x^2 + y^2). \label{eq:surface_integral_result}
\end{align}
\end{example}

% Source: Corsin Nick/Class Notes/Week 8.pdf, pp. 8--10
\begin{example}[A domain that Fubini alone cannot handle]
\label{ex:change_of_variables_domain}
Consider
\[\int_A \frac{y}{\sqrt{x}}\,dy\,dx, \qquad
A := \left\{(x,y) \in \mathbb{R}^2 \ :\ x,y>0,\ \ 1 \leq \frac{y}{\sqrt x} \leq 2,\ \ 1 \leq xy \leq 2\right\}.\]
This is not easily integrated with Fubini alone. But we can apply a transformation of variables
to simplify the integration domain $A$: take
\[\Phi : U \to V, \qquad \begin{pmatrix}x\\y\end{pmatrix} \mapsto
  \begin{pmatrix}u(x,y)\\v(x,y)\end{pmatrix} = \begin{pmatrix}\tfrac{y}{\sqrt x}\\[2pt] xy\end{pmatrix},\]
with $U$ and $V$ still to be determined, subject only to their being open.

\textbf{Check injectivity.} Suppose, for $x,y,x',y' > 0$,
\[\text{\textbf{(1)}}\ \ \frac{y}{\sqrt x} = \frac{y'}{\sqrt{x'}},
\qquad\qquad \text{\textbf{(2)}}\ \ xy = x'y'.\]
From \textbf{(2)}, $\dfrac{x}{x'} = \dfrac{y'}{y}$. Inserting into \textbf{(1)},
\[\frac{1}{\sqrt x} = \frac{y'}{y}\cdot\frac{1}{\sqrt{x'}} = \frac{x}{x'}\cdot\frac{1}{\sqrt{x'}}
\quad\Longrightarrow\quad x^{3/2} = (x')^{3/2} \quad\Longrightarrow\quad x = x',\]
and consequently $\dfrac{y'}{y} = \dfrac{x}{x'} = 1$. So, for $U := (0,\infty)^2$ and
$V := \Phi(U)$, the map $\Phi$ is a bijection between open sets.

\textbf{Functional determinant.} We calculate $|\det \Jac\Phi|$; if it is strictly positive we
may proceed, since $\Phi$ is then a local diffeomorphism as well as a bijection, hence a
diffeomorphism. Now
\[\Jac\Phi(x,y) = \begin{pmatrix} -\dfrac{y}{2\sqrt{x}^{\,3}} & \dfrac{1}{\sqrt x} \\[2ex] y & x \end{pmatrix},\]
so that
\[\big|\det \Jac\Phi(x,y)\big| = \left|-\frac32\,\frac{y}{\sqrt x}\right| = \frac32\,\frac{y}{\sqrt x} = \frac32\,u(x,y).\]
So, $\Phi$ is a diffeomorphism and $du\,dv = \tfrac32 u\,dx\,dy$. Furthermore, by design,
\[\Phi(A) = (1,2)^2.\]
So, with the change of variables (\cref{thm:change_of_variables}), followed by Fubini
(\cref{thm:fubini}),
\[\int_A \frac{y}{\sqrt x}\,dy\,dx
  = \int_{(1,2)^2} u\,\frac{1}{\tfrac32 u}\,du\,dv
  = \frac23\int_1^2 dv \int_1^2 du
  = \frac23.\]
\end{example}

\begin{ainote}
The design principle is worth stating, because it is the reverse of how substitution is usually
taught. Here $\Phi$ was not chosen to simplify the \emph{integrand} but to straighten the
\emph{domain}: the two constraints defining $A$ are literally $1 \leq u \leq 2$ and
$1 \leq v \leq 2$, so $\Phi$ turns an awkward curvilinear region into a square, and only then
does Fubini become available. That the integrand $y/\sqrt x$ is itself $u$, and cancels against
most of the Jacobian factor, is a bonus of the same choice.
\end{ainote}

\begin{aiexercise}[An order of integration that cannot be done]
% Generator: Claude Opus 5 (High)
\label{ex:ai_order_swap_gaussian}
Consider
\[I := \int_0^1\!\!\int_x^1 e^{y^2}\,dy\,dx.\]
\begin{enumerate}[label=\textbf{(\alph*)}]
  \item Explain what goes wrong if one attempts the inner integral as written.
  \item Describe the region of integration by inequalities, and rewrite $I$ with the two
    integrations in the opposite order.
  \item Evaluate $I$.
\end{enumerate}
\end{aiexercise}

% Source: exercises_2024/mock.pdf (Analysis II mock examination, Spring Semester 2024,
% Prof. Joaquim Serra), Exercise 12, p. 3
% Extractor: Claude Opus 5 (High)
% Filed with Fubini rather than with the improper integrals of chapter 19 because the whole
% problem is the product splitting; the one-dimensional pieces are then immediate. Part (a) is the
% trap and is what earns the block its place: the printed value sqrt(pi) is the answer to a
% one-dimensional Gaussian, and alpha = 0 is exactly the case where the z-integral does not exist.
\begin{exercise}[True or False \textnormal{(a Gaussian slab, and what happens at $\alpha=0$)}]
\label{ex:gaussian_slab_alpha_asymptotics}
For $\alpha \geq 0$ consider
\[I_\alpha := \int_{\mathbb{R}^3} e^{-x^2-y^2-\alpha\lvert z\rvert}\,dx\,dy\,dz .\]
Decide whether each of the following is true or false, with justification.
\begin{enumerate}[label=\textbf{(\alph*)}]
  \item $I_0 = \sqrt{\pi}$.
  \item $\alpha I_\alpha = I_1$.
  \item $\displaystyle\lim_{\alpha\to+\infty}\big(\alpha+\sqrt{\alpha}\big)I_\alpha = 2\pi$.
\end{enumerate}
\exinfo{This exercise is taken from the mock examination for Analysis II, Spring Semester 2024
(Prof.\ Joaquim Serra), where it appears as the twelfth block of multiple-choice questions. Its
paper supplies $\int_{\mathbb{R}}e^{-t^2}\,dt = \sqrt{\pi}$ as a formula that may be used without
proof; \cref{ex:gaussian_integral_polar} derives it.}
\end{exercise}

% Sonnet 5 (Medium)
We now want to look at a powerful new integration technique using ODEs.
